%=============================================================================
%  SBR 2026 - Brazilian Symposium on Robotics
%  Joao Pessoa, PB, Brazil - November 24-27, 2026
%  https://www.natalnet.br/sbr2026/
%
%  VERSAO DE SUBMISSAO - DOUBLE BLIND
%  ---------------------------------------------------------------------------
%  REGRAS OBRIGATORIAS (nao violar):
%   * Ingles.
%   * Maximo 6 paginas, INCLUINDO figuras, tabelas e referencias.
%   * Classe IEEEtran, opcao "conference". NAO alterar margens/fontes.
%   * NENHUMA identificacao de autor, afiliacao, e-mail, cidade, laboratorio,
%     agencia de fomento, numero de processo ou agradecimento.
%   * Autocitacoes SEMPRE em terceira pessoa.
%   * Submissao: https://jems3.sbc.org.br/sbr2026  (deadline 30/07/2026)
%
%  NA CAMERA-READY (15/09/2026): reinserir autores + Acknowledgment
%  (procure pelos blocos marcados com "CAMERA-READY").
%=============================================================================

\documentclass[conference]{IEEEtran}
\IEEEoverridecommandlockouts
% Se o compilador reclamar de \IEEEoverridecommandlockouts, apenas comente a linha.

%----------------------------- Pacotes ---------------------------------------
\usepackage[T1]{fontenc}
\usepackage[utf8]{inputenc}
\usepackage{amsmath,amssymb,amsfonts}
\usepackage{graphicx}
\usepackage{booktabs}
\usepackage{subcaption}       % NAO use "subfigure" (obsoleto)
\usepackage{algorithm}
\usepackage{algpseudocode}
\usepackage{textcomp}
\usepackage{xcolor}
\usepackage{url}
\usepackage{cite}
\usepackage[hidelinks]{hyperref}   % sempre por ultimo (antes de cleveref)
\usepackage{cleveref}

%--------------------- Limpeza de metadados do PDF ---------------------------
% CRITICO PARA O DOUBLE-BLIND: impede que o nome do autor vaze no PDF.
\hypersetup{
  pdftitle={SBR 2026 Submission},
  pdfauthor={},
  pdfsubject={},
  pdfkeywords={},
  pdfcreator={}
}

\graphicspath{{fig/}}

%--------------------------- Macros uteis ------------------------------------
\newcommand{\ours}{the proposed method}   % evita "our method" em excesso

% Placeholder de figura: permite compilar o esqueleto SEM os arquivos de imagem.
% Ao inserir as figuras reais, troque \figph{...} por \includegraphics{...}.
\newcommand{\figph}[2][\columnwidth]{%
  \fbox{\parbox[c][2.6cm][c]{\dimexpr#1-2\fboxsep-2\fboxrule\relax}{%
    \centering\footnotesize\texttt{#2}\\[2pt]
    (figure placeholder --- replace with \texttt{\textbackslash includegraphics})}}}
\DeclareMathOperator*{\argmin}{arg\,min}
\DeclareMathOperator*{\argmax}{arg\,max}

\begin{document}

%============================== TITULO =======================================
% Titulo: informativo, sem siglas obscuras, sem ":" duplo.
\title{Title of the Paper: Concise, Specific and Informative\\
{\footnotesize \textit{(remove any subtitle that identifies a project or lab)}}}

%============================== AUTORES ======================================
% >>> VERSAO DE SUBMISSAO (DOUBLE BLIND) - MANTENHA ESTE BLOCO <<<
\author{\IEEEauthorblockN{Anonymous Author(s)}
\IEEEauthorblockA{\textit{Affiliation omitted for double-blind review}\\
\textit{Institution omitted}\\
Contact information omitted}}

% >>> CAMERA-READY: apague o bloco acima e descomente o modelo abaixo <<<
% \author{
% \IEEEauthorblockN{1\textsuperscript{st} Given Name Surname}
% \IEEEauthorblockA{\textit{dept. name of organization} \\
% \textit{name of organization}\\ City, Country \\ email@domain}
% \and
% \IEEEauthorblockN{2\textsuperscript{nd} Given Name Surname}
% \IEEEauthorblockA{\textit{dept. name of organization} \\
% \textit{name of organization}\\ City, Country \\ email@domain}
% }

\maketitle

%============================== ABSTRACT =====================================
% 150-250 palavras. Sem citacoes, sem siglas nao definidas, sem equacoes.
% Formula: contexto -> lacuna -> o que foi feito -> como -> resultado NUMERICO -> impacto.
\begin{abstract}
One or two sentences of context on the robotics problem.
One sentence stating the specific gap or limitation of existing approaches.
This paper presents \emph{X}, a \ldots{} that \ldots.
The method combines \ldots{} with \ldots, and is evaluated on \ldots{}
(simulation in Gazebo / dataset \ldots{} / a real differential-drive platform).
Experiments against \emph{N} baselines show a reduction of \emph{YY}\,\% in
\emph{metric}, while running at \emph{Z}\,Hz on \emph{hardware class}.
The results indicate that \ldots.
\end{abstract}

%============================== KEYWORDS =====================================
% 3-6 termos, alinhados aos topicos de interesse do SBR 2026.
\begin{IEEEkeywords}
mobile robotics, motion planning, deep learning, autonomous navigation, ROS~2
\end{IEEEkeywords}

%============================ I. INTRODUCTION ================================
\section{Introduction}
\label{sec:intro}

Paragraph 1 --- \textbf{Context}: why this robotics problem matters
(application, safety, autonomy level).

Paragraph 2 --- \textbf{State of the art in one breath}: classical approaches
\cite{ref_classic} and learning-based approaches \cite{ref_learning} have
addressed \ldots.

Paragraph 3 --- \textbf{Gap}: however, existing methods
(i) assume \ldots, (ii) do not scale to \ldots, and (iii) were not validated
under \ldots.

Paragraph 4 --- \textbf{Proposal}: this paper proposes \ldots.

\noindent The main contributions of this work are:
\begin{enumerate}
  \item A \ldots{} formulation that \ldots;
  \item A \ldots{} architecture/algorithm that \ldots;
  \item An experimental evaluation on \ldots{} showing \ldots{} against
        \emph{N} baselines, including an ablation study of \ldots.
\end{enumerate}

The remainder of this paper is organized as follows.
Section~\ref{sec:related} reviews related work.
Section~\ref{sec:method} describes the proposed approach.
Section~\ref{sec:setup} presents the experimental setup and
Section~\ref{sec:results} discusses the results.
Section~\ref{sec:conclusion} concludes the paper.

%============================ II. RELATED WORK ===============================
\section{Related Work}
\label{sec:related}

% Organize por FAMILIA de abordagem, nao por artigo ("[1] did X. [2] did Y.").
% ATENCAO DOUBLE-BLIND: cite trabalhos proprios em TERCEIRA PESSOA.
%   ERRADO: "In our previous work [7], we showed..."
%   CERTO : "Silva et al. [7] proposed ..., which is adopted here as baseline."

\subsection{Classical approaches}
Methods based on \ldots{} \cite{ref_classic} \ldots.

\subsection{Learning-based approaches}
More recently, \ldots{} \cite{ref_learning,ref_transformer} \ldots.

\subsection{Positioning of this work}
Unlike the works above, \ours{} \ldots. To the best of the authors'
knowledge, this is the first work that \ldots.
% Note: "to the best of the authors' knowledge" (nao "our knowledge") mantem
% o tom impessoal, embora "our" seja aceitavel em double-blind desde que nao
% revele identidade. Prefira a forma impessoal.

%============================ III. METHODOLOGY ===============================
\section{Proposed Approach}
\label{sec:method}

\subsection{Problem formulation}
Let $\mathcal{X} \subset \mathbb{R}^n$ denote \ldots{} and let
$x_t$ be the state at time $t$. The goal is to find

\begin{equation}
  \pi^{\star} = \argmin_{\pi \in \Pi} \; \mathbb{E}\!\left[ \sum_{t=0}^{T} c(x_t, \pi(x_t)) \right],
  \label{eq:problem}
\end{equation}

\noindent subject to \ldots. Every symbol used in \eqref{eq:problem} must be
defined at first use.

\subsection{System overview}
Figure~\ref{fig:arch} shows the overall architecture, composed of
(i) \ldots, (ii) \ldots{} and (iii) \ldots.

\begin{figure}[!t]
  \centering
  % Use PDF vetorial para diagramas; PNG >= 300 dpi para fotos.
  % ATENCAO: remova logos institucionais e nomes de pessoas das imagens.
  % \includegraphics[width=\columnwidth]{arch.pdf}
  \figph{fig/arch.pdf}
  \caption{Overview of the proposed architecture. The perception module
  produces \ldots, which feeds the \ldots{} module.}
  \label{fig:arch}
\end{figure}

\subsection{Module A}
\ldots

\subsection{Module B}
\ldots

\begin{algorithm}[t]
\caption{Proposed procedure}
\label{alg:main}
\begin{algorithmic}[1]
\Require initial state $x_0$, goal $g$, model $f_\theta$
\Ensure trajectory $\tau$
\State $\tau \gets \{x_0\}$
\For{$t = 0$ \textbf{to} $T-1$}
  \State $u_t \gets \pi_\theta(x_t, g)$
  \State $x_{t+1} \gets f(x_t, u_t)$
  \State $\tau \gets \tau \cup \{x_{t+1}\}$
\EndFor
\State \Return $\tau$
\end{algorithmic}
\end{algorithm}

%========================== IV. EXPERIMENTAL SETUP ===========================
\section{Experimental Setup}
\label{sec:setup}

\subsection{Datasets / environments}
% Se o dataset for do proprio grupo, descreva-o sem nomear o laboratorio:
% "an in-house dataset collected with a differential-drive platform".
\ldots

\subsection{Baselines}
\ours{} is compared against \cite{ref_classic}, \cite{ref_learning} and
\cite{ref_transformer}, all reimplemented/retrained under identical conditions.

\subsection{Metrics}
\ldots{} (define each metric with its unit and whether lower/higher is better).

\subsection{Implementation details}
% Reprodutibilidade: versoes, hardware, hiperparametros, tempo de treino.
Experiments were run on \emph{hardware class} (e.g., an NVIDIA RTX-class GPU),
using ROS~2 Humble and PyTorch~2.x. Training used \ldots{} for \ldots{} epochs
with learning rate \ldots. All results are averaged over \emph{k} runs with
different random seeds; mean and standard deviation are reported.
% DOUBLE-BLIND: nada de "our cluster at <universidade>".

% Codigo/dados: anonimize ou omita.
The source code will be made publicly available upon acceptance
(link omitted for double-blind review).

%======================== V. RESULTS AND DISCUSSION ==========================
\section{Results and Discussion}
\label{sec:results}

\subsection{Quantitative results}
Table~\ref{tab:results} reports \ldots.

\begin{table}[!t]
  \caption{Comparison against baselines. Best results in bold.
           $\downarrow$ indicates that lower is better.}
  \label{tab:results}
  \centering
  \begin{tabular}{lccc}
    \toprule
    Method & Metric A $\downarrow$ & Metric B $\downarrow$ & Runtime (ms) \\
    \midrule
    Baseline~1 \cite{ref_classic}      & 0.52 & 1.10 & \textbf{8} \\
    Baseline~2 \cite{ref_learning}     & 0.47 & 0.99 & 21 \\
    Baseline~3 \cite{ref_transformer}  & 0.44 & 0.93 & 34 \\
    \textbf{Proposed}                  & \textbf{0.41} & \textbf{0.88} & 15 \\
    \bottomrule
  \end{tabular}
\end{table}

\subsection{Ablation study}
Table~\ref{tab:ablation} isolates the contribution of each component.

\begin{table}[!t]
  \caption{Ablation study of the proposed components.}
  \label{tab:ablation}
  \centering
  \begin{tabular}{ccc c}
    \toprule
    Comp.~A & Comp.~B & Comp.~C & Metric A $\downarrow$ \\
    \midrule
    \checkmark &            &            & 0.50 \\
    \checkmark & \checkmark &            & 0.45 \\
    \checkmark & \checkmark & \checkmark & \textbf{0.41} \\
    \bottomrule
  \end{tabular}
\end{table}

\subsection{Qualitative results}
Figure~\ref{fig:qualitative} illustrates \ldots.

\begin{figure*}[!t]   % ocupa as duas colunas
  \centering
  % Nas versoes finais, troque \figph por:
  %   \includegraphics[width=\textwidth]{qual_a.pdf}
  \begin{subfigure}[b]{0.32\textwidth}
    \figph[\textwidth]{fig/qual\_a.pdf}\caption{Scenario A}\label{fig:qa}
  \end{subfigure}\hfill
  \begin{subfigure}[b]{0.32\textwidth}
    \figph[\textwidth]{fig/qual\_b.pdf}\caption{Scenario B}\label{fig:qb}
  \end{subfigure}\hfill
  \begin{subfigure}[b]{0.32\textwidth}
    \figph[\textwidth]{fig/qual\_c.pdf}\caption{Scenario C}\label{fig:qc}
  \end{subfigure}
  \caption{Qualitative comparison in three representative scenarios.
  Curves are distinguishable in grayscale.}
  \label{fig:qualitative}
\end{figure*}

\subsection{Discussion and limitations}
The results indicate that \ldots. Nevertheless, \ours{} presents the following
limitations: (i) \ldots; (ii) \ldots. These aspects motivate the extensions
discussed in Section~\ref{sec:conclusion}.

%============================ VI. CONCLUSION =================================
\section{Conclusion and Future Work}
\label{sec:conclusion}

This paper presented \ldots. Experiments on \ldots{} showed \ldots.
Future work includes \ldots, \ldots{} and validation on a physical platform
under \ldots.
% Nao apresente resultados novos na conclusao.

%========================= ACKNOWLEDGMENT (CAMERA-READY) =====================
% !!! NAO INCLUIR NA SUBMISSAO DOUBLE-BLIND !!!
% Descomente somente na camera-ready:
%
% \section*{Acknowledgment}
% The authors thank <agencia> for financial support under grant <numero>.

%============================== REFERENCIAS ==================================
\bibliographystyle{IEEEtran}
\bibliography{references}

% Alternativa sem BibTeX (se preferir lista manual):
% \begin{thebibliography}{00}
% \bibitem{ref_classic} S. Thrun, W. Burgard, and D. Fox, \emph{Probabilistic Robotics}. Cambridge, MA, USA: MIT Press, 2005.
% \end{thebibliography}

\end{document}
